% ═══════════════════════════════════════════════════════════════════════════════
% Chapter 10: TetraData Structure and TetraMemory
% Part III: The DOGMA Architecture
% ═══════════════════════════════════════════════════════════════════════════════

\chapter{TetraData Structure and TetraMemory}
\label{ch:tetradata}

\chapterepigraph{The tetrahedron is the simplest possible solid. Any geometry must begin here.}{After Buckminster Fuller}

\noindent
Computer graphics has long relied on \emph{triangles} as the fundamental rendering primitive.
Triangle meshes approximate surfaces efficiently and map naturally to GPU hardware.
But triangles are inherently \emph{two-dimensional}: they describe surfaces, not volumes.
For a computational architecture inspired by three-dimensional molecular biology---where
proteins fold into globular structures, DNA coils into superhelical conformations, and
tetrahedral carbon bonds define organic chemistry---a volumetric primitive is needed.

This chapter introduces the \textbf{TetraData Structure} (TDS): a computational framework
built on \emph{tetrahedra} rather than triangles. We begin with a patient, step-by-step
construction that builds a tetrahedron one vertex at a time, then show how tetrahedra
link together into chains that map directly onto DNA sequences. We then develop
\textbf{TetraMemory}, a structured memory system that replaces self-attention in DOGMA
with local tetrahedral interactions that compose into global computation.

% ═══════════════════════════════════════════════════════════════════════════════
\section{Building a Tetrahedron from Scratch}
\label{sec:building-tetra}
% ═══════════════════════════════════════════════════════════════════════════════

Before defining the TetraData Structure formally, let us build our fundamental
primitive---the tetrahedron---one vertex at a time. This construction will reveal
\emph{why} the tetrahedron is the minimal solid that carries volume and chirality,
and why no simpler primitive will do.

% ─── Step 0: A Single Vertex ─────────────────────────────────────────────────
\subsection{Step 0: A Single Vertex (0D)}
\label{sec:step0-vertex}

We begin with nothing: an empty space. We place a single point $v_0$ and assign
it a value (say, a DNA base, a token embedding, or simply a label).

\begin{center}
\begin{tikzpicture}[scale=1.2]
  % Single vertex
  \fill[dogmablue] (0,0) circle (4pt);
  \node[above=6pt, font=\bfseries\small, color=dogmablue] at (0,0) {$v_0$};

  % Annotations
  \node[below=20pt, font=\small, align=center] at (0,0) {
    Vertices: 1 \quad Edges: 0 \quad Faces: 0 \quad Volume: 0
  };

  % Dimension label
  \node[right=30pt, font=\itshape\small, color=dogmagray] at (0,0) {0-dimensional};
\end{tikzpicture}
\end{center}

A single vertex is a 0-dimensional object. It can store a value---but it encodes
no \emph{relationships}. There is no notion of distance, connection, area, or volume.
In a sequence model, a single vertex is a single token in isolation, with no
context whatsoever.

\begin{remark}
In graph theory terms, we have the graph $K_1$: one vertex, zero edges. The
adjacency matrix is the $1 \times 1$ zero matrix. This is the least informative
data structure possible.
\end{remark}

% ─── Step 1: Two Vertices ────────────────────────────────────────────────────
\subsection{Step 1: Add a Second Vertex (1D --- an Edge)}
\label{sec:step1-edge}

Now add a second vertex $v_1$ and connect it to $v_0$ with an edge $e_{01}$.

\begin{center}
\begin{tikzpicture}[scale=1.2]
  % Two vertices with edge
  \fill[dogmablue] (0,0) circle (4pt);
  \fill[dogmablue] (3,0) circle (4pt);
  \draw[thick, dogmablue] (0,0) -- (3,0);

  \node[above=6pt, font=\bfseries\small, color=dogmablue] at (0,0) {$v_0$};
  \node[above=6pt, font=\bfseries\small, color=dogmablue] at (3,0) {$v_1$};
  \node[below=4pt, font=\small, color=dogmateal] at (1.5,0) {$e_{01}$};

  % Annotations
  \node[below=20pt, font=\small, align=center] at (1.5,0) {
    Vertices: 2 \quad Edges: 1 \quad Faces: 0 \quad Volume: 0
  };

  \node[right=15pt, font=\itshape\small, color=dogmagray] at (3,0) {1-dimensional};
\end{tikzpicture}
\end{center}

An edge is a 1-dimensional object. It encodes exactly one relationship: the
pairwise interaction between $v_0$ and $v_1$. The edge can carry a weight
$w_{01} \in \mathbb{R}$ representing the strength or character of that interaction.

In the language of neural networks, this is a single attention weight: how much
does token $v_0$ attend to token $v_1$? One edge captures one pairwise interaction.
This is the graph $K_2$: the complete graph on 2 vertices.

\begin{remark}
An edge has no \emph{area}, no \emph{volume}, and no \emph{handedness}. You
cannot distinguish a ``left-handed'' line segment from a ``right-handed'' one.
An edge is symmetric under reflection.
\end{remark}

% ─── Step 2: Three Vertices ──────────────────────────────────────────────────
\subsection{Step 2: Add a Third Vertex (2D --- a Triangle)}
\label{sec:step2-triangle}

Add a third vertex $v_2$ and connect it to both existing vertices. We now have
three vertices and three edges forming a triangle.

\begin{center}
\begin{tikzpicture}[scale=1.4]
  % Triangle
  \coordinate (v0) at (0,0);
  \coordinate (v1) at (3,0);
  \coordinate (v2) at (1.5,2.2);

  % Fill
  \fill[dogmalight, opacity=0.5] (v0) -- (v1) -- (v2) -- cycle;

  % Edges
  \draw[thick, dogmablue] (v0) -- (v1);
  \draw[thick, dogmablue] (v1) -- (v2);
  \draw[thick, dogmablue] (v2) -- (v0);

  % Vertices
  \fill[dogmablue] (v0) circle (4pt);
  \fill[dogmablue] (v1) circle (4pt);
  \fill[dogmablue] (v2) circle (4pt);

  % Labels
  \node[below left=3pt, font=\bfseries\small, color=dogmablue] at (v0) {$v_0$};
  \node[below right=3pt, font=\bfseries\small, color=dogmablue] at (v1) {$v_1$};
  \node[above=5pt, font=\bfseries\small, color=dogmablue] at (v2) {$v_2$};

  % Edge labels
  \node[below=3pt, font=\small, color=dogmateal] at (1.5,0) {$e_{01}$};
  \node[right=3pt, font=\small, color=dogmateal] at ($(v1)!0.5!(v2)$) {$e_{12}$};
  \node[left=3pt, font=\small, color=dogmateal] at ($(v0)!0.5!(v2)$) {$e_{02}$};

  % Annotations
  \node[below=25pt, font=\small, align=center] at (1.5,0) {
    Vertices: 3 \quad Edges: 3 \quad Faces: 1 \quad Volume: 0
  };

  \node[right=15pt, font=\itshape\small, color=dogmagray] at (v1) {2-dimensional};
\end{tikzpicture}
\end{center}

A triangle is a 2-dimensional object. It has \emph{area} but no \emph{volume}.
Three edges encode all $\binom{3}{2} = 3$ pairwise interactions among the three
vertices, forming the complete graph $K_3$.

This is the primitive used in computer graphics: triangle meshes tile surfaces.
GPUs are optimized for rasterizing triangles. But triangles have two critical
limitations for our purposes:

\begin{enumerate}[leftmargin=2em]
  \item \textbf{No volume.} A triangle is flat. It cannot enclose space.
        It describes a boundary, not an interior.
  \item \textbf{No chirality.} A triangle in 2D can be reflected into its
        mirror image by a rotation. In 3D, a triangle embedded in a plane
        does have two ``sides,'' but this orientation is not an intrinsic
        property---it depends on the choice of normal direction.
\end{enumerate}

\begin{bioinsight}[Why Triangles Are Not Enough]
In molecular biology, chirality is not optional---it is essential.
All naturally occurring amino acids are L-amino acids (left-handed).
DNA's double helix is always right-handed (B-form).
Enzymes distinguish between mirror-image molecules with exquisite specificity.
A computational primitive that cannot represent chirality is fundamentally
incapable of encoding molecular structure faithfully.
A triangle mesh can approximate the \emph{surface} of a protein, but it
cannot capture the \emph{handed} internal packing that determines function.
\end{bioinsight}

% ─── Step 3: Four Vertices ───────────────────────────────────────────────────
\subsection{Step 3: Add a Fourth Vertex (3D --- the Tetrahedron!)}
\label{sec:step3-tetrahedron}

Now the crucial step. Add a fourth vertex $v_3$ that is \emph{not in the plane}
of the first three, and connect it to all existing vertices. We obtain a
tetrahedron: the simplest possible solid.

\begin{center}
\begin{tikzpicture}[scale=1.6]
  % Vertices of the tetrahedron (3D projection)
  \coordinate (v0) at (0,0);
  \coordinate (v1) at (3.5,0);
  \coordinate (v2) at (1.75,1.8);
  \coordinate (v3) at (1.2,0.7);  % interior-ish vertex for 3D effect

  % Back faces (lighter)
  \fill[dogmamint, opacity=0.3] (v0) -- (v1) -- (v3) -- cycle;
  \fill[dogmalight, opacity=0.3] (v0) -- (v2) -- (v3) -- cycle;

  % Front face
  \fill[dogmalight, opacity=0.5] (v0) -- (v1) -- (v2) -- cycle;
  \fill[dogmamint, opacity=0.4] (v1) -- (v2) -- (v3) -- cycle;

  % Back edges (dashed)
  \draw[thick, dogmateal, dashed] (v0) -- (v3);
  \draw[thick, dogmateal, dashed] (v1) -- (v3);

  % Front edges (solid)
  \draw[thick, dogmablue] (v0) -- (v1);
  \draw[thick, dogmablue] (v1) -- (v2);
  \draw[thick, dogmablue] (v2) -- (v0);
  \draw[thick, dogmablue] (v2) -- (v3);

  % Vertices
  \fill[dogmablue] (v0) circle (4pt);
  \fill[dogmablue] (v1) circle (4pt);
  \fill[dogmablue] (v2) circle (4pt);
  \fill[dogmateal] (v3) circle (4pt);

  % Labels
  \node[below left=3pt, font=\bfseries\small, color=dogmablue] at (v0) {$v_0$};
  \node[below right=3pt, font=\bfseries\small, color=dogmablue] at (v1) {$v_1$};
  \node[above=5pt, font=\bfseries\small, color=dogmablue] at (v2) {$v_2$};
  \node[below=5pt, font=\bfseries\small, color=dogmateal] at (v3) {$v_3$};

  % Edge count annotation
  \node[below=25pt, font=\small, align=center] at (1.75,0) {
    Vertices: 4 \quad Edges: 6 \quad Faces: 4 \quad Volume: $>0$
  };

  \node[right=10pt, font=\itshape\small, color=dogmagray] at (v1) {3-dimensional!};

  % Chirality indicator
  \draw[-{Stealth}, thick, dogmared, opacity=0.6] (1.6,0.9) arc[start angle=180, end angle=60, radius=0.3];
  \node[right=2pt, font=\scriptsize, color=dogmared] at (1.9,1.1) {$\chi$};
\end{tikzpicture}
\end{center}

A tetrahedron is the \emph{first} object in our construction that possesses:
\begin{itemize}[leftmargin=2em]
  \item \textbf{Volume:} It encloses a region of 3D space. The signed volume is
        computed from the $3 \times 3$ determinant of the edge vectors.
  \item \textbf{Chirality:} The sign of the determinant ($+1$ or $-1$) defines
        a handedness that cannot be changed by any rotation---only by reflection.
  \item \textbf{Complete pairwise interactions:} Six edges encode all
        $\binom{4}{2} = 6$ pairwise relationships. This is the complete graph $K_4$.
  \item \textbf{Four faces:} Each face is a triangle that can serve as an
        interface for connecting to another tetrahedron.
\end{itemize}

\begin{definition}[Tetrahedron]
\label{def:tetradata-tetrahedron}
A \emph{tetrahedron} $\mathcal{T}$ in $\mathbb{R}^3$ is the convex hull of four non-coplanar points $\mathbf{v}_0, \mathbf{v}_1, \mathbf{v}_2, \mathbf{v}_3$:
\begin{equation}
\mathcal{T} = \left\{ \sum_{i=0}^{3} \lambda_i \mathbf{v}_i \;\middle|\; \lambda_i \geq 0, \; \sum_{i=0}^{3} \lambda_i = 1 \right\}.
\end{equation}
The \emph{signed volume} is:
\begin{equation}
V_{\text{signed}}(\mathcal{T}) = \frac{1}{6} \det \begin{pmatrix} \mathbf{v}_1 - \mathbf{v}_0 & \mathbf{v}_2 - \mathbf{v}_0 & \mathbf{v}_3 - \mathbf{v}_0 \end{pmatrix},
\label{eq:signed-volume}
\end{equation}
and the \emph{chirality} is:
\begin{equation}
\chi(\mathcal{T}) = \operatorname{sign}\!\big(V_{\text{signed}}(\mathcal{T})\big) \in \{+1, -1\}.
\end{equation}
\end{definition}

\begin{keyinsight}[The Tetrahedron Is the Minimal Solid]
In all of Euclidean geometry, the tetrahedron is the \emph{simplest} object that
simultaneously possesses volume and chirality. You cannot build a solid with
fewer than 4 vertices. This is not a design choice---it is a mathematical
necessity. Just as the triangle is the minimal polygon, the tetrahedron is
the minimal polyhedron.

DOGMA's choice of the tetrahedron as its fundamental primitive is therefore
not arbitrary: it is the minimal structure that can encode 3D geometry,
handedness, and complete pairwise interactions among a set of elements.
\end{keyinsight}

% ─── Step-by-step summary figure ──────────────────────────────────────────────
\subsection{The Construction at a Glance}

The following figure summarizes the four steps side by side, highlighting what
each new vertex contributes:

\begin{figure}[htbp]
\centering
\begin{tikzpicture}[scale=0.95]
  % ── Step 0: single vertex ──
  \begin{scope}[shift={(0,0)}]
    \fill[dogmablue] (0,0) circle (3pt);
    \node[above=5pt, font=\scriptsize\bfseries] at (0,0) {$v_0$};
    \node[below=12pt, font=\scriptsize, align=center] at (0,0) {
      \textbf{Step 0} \\ $V{=}1$, $E{=}0$ \\ 0D
    };
  \end{scope}

  % Arrow
  \draw[-{Stealth}, thick, dogmagray] (0.6,0) -- (1.8,0);

  % ── Step 1: edge ──
  \begin{scope}[shift={(2.5,0)}]
    \fill[dogmablue] (0,0) circle (3pt);
    \fill[dogmablue] (1.2,0) circle (3pt);
    \draw[thick, dogmablue] (0,0) -- (1.2,0);
    \node[above=5pt, font=\scriptsize\bfseries] at (0,0) {$v_0$};
    \node[above=5pt, font=\scriptsize\bfseries] at (1.2,0) {$v_1$};
    \node[below=12pt, font=\scriptsize, align=center] at (0.6,0) {
      \textbf{Step 1} \\ $V{=}2$, $E{=}1$ \\ 1D
    };
  \end{scope}

  % Arrow
  \draw[-{Stealth}, thick, dogmagray] (4.3,0) -- (5.5,0);

  % ── Step 2: triangle ──
  \begin{scope}[shift={(6.2,0)}]
    \coordinate (a) at (0,0);
    \coordinate (b) at (1.4,0);
    \coordinate (c) at (0.7,1.0);
    \fill[dogmalight, opacity=0.5] (a) -- (b) -- (c) -- cycle;
    \draw[thick, dogmablue] (a) -- (b) -- (c) -- cycle;
    \fill[dogmablue] (a) circle (3pt);
    \fill[dogmablue] (b) circle (3pt);
    \fill[dogmablue] (c) circle (3pt);
    \node[below left=1pt, font=\scriptsize\bfseries] at (a) {$v_0$};
    \node[below right=1pt, font=\scriptsize\bfseries] at (b) {$v_1$};
    \node[above=3pt, font=\scriptsize\bfseries] at (c) {$v_2$};
    \node[below=12pt, font=\scriptsize, align=center] at (0.7,0) {
      \textbf{Step 2} \\ $V{=}3$, $E{=}3$ \\ 2D (triangle)
    };
  \end{scope}

  % Arrow
  \draw[-{Stealth}, thick, dogmagray] (8.2,0) -- (9.4,0);

  % ── Step 3: tetrahedron ──
  \begin{scope}[shift={(10.1,0)}]
    \coordinate (a) at (0,0);
    \coordinate (b) at (1.6,0);
    \coordinate (c) at (0.8,1.1);
    \coordinate (d) at (0.5,0.4);
    \fill[dogmamint, opacity=0.3] (a) -- (b) -- (d) -- cycle;
    \fill[dogmalight, opacity=0.5] (a) -- (b) -- (c) -- cycle;
    \fill[dogmamint, opacity=0.4] (b) -- (c) -- (d) -- cycle;
    \draw[thick, dogmablue] (a) -- (b) -- (c) -- cycle;
    \draw[thick, dogmateal, dashed] (a) -- (d);
    \draw[thick, dogmateal, dashed] (b) -- (d);
    \draw[thick, dogmablue] (c) -- (d);
    \fill[dogmablue] (a) circle (3pt);
    \fill[dogmablue] (b) circle (3pt);
    \fill[dogmablue] (c) circle (3pt);
    \fill[dogmateal] (d) circle (3pt);
    \node[below left=1pt, font=\scriptsize\bfseries] at (a) {$v_0$};
    \node[below right=1pt, font=\scriptsize\bfseries] at (b) {$v_1$};
    \node[above=3pt, font=\scriptsize\bfseries] at (c) {$v_2$};
    \node[below=2pt, font=\scriptsize\bfseries, color=dogmateal] at (d) {$v_3$};
    \node[below=12pt, font=\scriptsize, align=center] at (0.8,0) {
      \textbf{Step 3} \\ $V{=}4$, $E{=}6$ \\ 3D (tetrahedron!)
    };
  \end{scope}
\end{tikzpicture}
\caption{Building a tetrahedron one vertex at a time. Each step adds a new
dimension: a vertex (0D), an edge (1D), a triangle (2D), and finally a
tetrahedron (3D). The tetrahedron is the first object with volume and chirality.}
\label{fig:tetra-construction}
\end{figure}

\begin{example}[Counting Combinatorics]
At each step $k$ (adding vertex $v_k$ to the existing structure), the new
vertex connects to all $k$ existing vertices, adding exactly $k$ new edges:
\begin{center}
\renewcommand{\arraystretch}{1.2}
\begin{tabular}{ccccl}
\toprule
\textbf{Step} & \textbf{Vertices} & \textbf{New edges} & \textbf{Total edges} & \textbf{Graph} \\
\midrule
0 & 1 & 0 & 0 & $K_1$ \\
1 & 2 & 1 & 1 & $K_2$ \\
2 & 3 & 2 & 3 & $K_3$ \\
3 & 4 & 3 & 6 & $K_4$ \\
\bottomrule
\end{tabular}
\end{center}
In general, the complete graph $K_n$ has $\binom{n}{2} = \frac{n(n-1)}{2}$ edges.
For $K_4$, this gives $\binom{4}{2} = 6$ edges---exactly the six edges of a
tetrahedron.
\end{example}

% ═══════════════════════════════════════════════════════════════════════════════
\section{Why Tetrahedra? --- A Deeper Look}
\label{sec:why-tetra}
% ═══════════════════════════════════════════════════════════════════════════════

With the construction complete, we can appreciate why DOGMA chooses the
tetrahedron over simpler primitives.

\subsection{Tetrahedra vs.\ Triangles: A Detailed Comparison}

\begin{center}
\renewcommand{\arraystretch}{1.15}
\begin{tabularx}{0.95\textwidth}{L{3cm} L{4cm} Y}
\toprule
\textbf{Property} & \textbf{Triangle Mesh} & \textbf{Tetrahedral Mesh (TDS)} \\
\midrule
Dimension & 2D surface in 3D & 3D volume \\
Vertices per element & 3 & 4 \\
Edges per element & 3 ($K_3$) & 6 ($K_4$) --- $2\times$ more interactions \\
Faces per element & 1 & 4 triangular faces \\
Interior & None (boundary only) & Enclosed volume \\
Chirality & None & $\pm 1$ handedness \\
Adjacency & Edge-sharing & Face-sharing (richer interface) \\
DNA mapping & --- & 4 vertices $\leftrightarrow$ \{A, C, G, T\} \\
DOF per element & 9 (positions only) & 36 (positions + edges + bases + vol + chir) \\
Info.\ density & 3 edges & 6 edges --- $4\times$ information density\\
\bottomrule
\end{tabularx}
\end{center}

The last row deserves emphasis. A triangle encodes 3 pairwise interactions.
A tetrahedron encodes 6. But the tetrahedron has only one more vertex.
In terms of pairwise interactions per vertex, the triangle achieves
$3/3 = 1$ edge per vertex, while the tetrahedron achieves $6/4 = 1.5$ edges per
vertex---a 50\% improvement in information density.

\subsection{The Biological Justification}

The biological case for tetrahedra is compelling:

\begin{itemize}[leftmargin=1.8em]
  \item \textbf{Carbon chemistry:} The carbon atom---the backbone of all organic
        molecules---forms four bonds arranged tetrahedrally ($sp^3$ hybridization).
        Methane (CH$_4$) is a perfect tetrahedron. Protein backbone geometry is
        fundamentally tetrahedral at each $C_\alpha$ center.

  \item \textbf{DNA base count:} DNA uses exactly four bases---A, C, G, T. A
        tetrahedron has exactly four vertices. This numerical coincidence becomes
        the foundation of DOGMA's data structure: each vertex of a tetrahedron
        is assigned one DNA base.

  \item \textbf{DNA base stacking:} Adjacent base pairs in the double helix create
        local geometries that are naturally described by tetrahedral coordinates.

  \item \textbf{Protein structure:} Alpha-carbon ($C_\alpha$) geometry in protein
        backbones forms tetrahedral arrangements that are native to TDS
        representation.

  \item \textbf{Chirality in biology:} All naturally occurring amino acids are
        L-amino acids; DNA forms a right-handed helix. Chirality is
        \emph{information} that must be preserved, and only volumetric
        primitives (3D and above) can encode it.
\end{itemize}

\begin{bioinsight}[Four Bases, Four Vertices]
The fact that DNA uses exactly four nucleotide bases is one of the most
fundamental constraints in molecular biology. This four-fold alphabet maps
perfectly onto the four vertices of a tetrahedron. This is not merely a
convenient analogy: it is the mathematical foundation of the TetraData
Structure. Each vertex carries a base type, and the six edges encode all
pairwise interactions among those bases---including the Watson--Crick pairings
(A--T and G--C) as geometrically distinguished edges.
\end{bioinsight}

% ═══════════════════════════════════════════════════════════════════════════════
\section{Linking Tetrahedra --- The LinkedTetraChain}
\label{sec:linked-tetra-chain}
% ═══════════════════════════════════════════════════════════════════════════════

A single tetrahedron encodes the interactions among four elements. But real
sequences are longer than four tokens. How do we handle sequences of arbitrary
length? The answer is \emph{linking}: we chain tetrahedra together by sharing
triangular faces, exactly as the PowerPoint construction shows.

\subsection{Adding Vertex 4: The Second Tetrahedron}
\label{sec:add-v4}

We continue the construction from Section~\ref{sec:building-tetra}. Having
built the first tetrahedron $\mathcal{T}_1 = (v_0, v_1, v_2, v_3)$, we now add
vertex $v_4$ and connect it to the three most recent vertices: $v_1, v_2, v_3$.
This creates a \emph{second} tetrahedron:
\[
  \mathcal{T}_2 = (v_1, v_2, v_3, v_4).
\]

The two tetrahedra share the triangular face $(v_1, v_2, v_3)$. The new vertex
$v_4$ adds exactly 3 new edges: $e_{14}, e_{24}, e_{34}$.

\begin{center}
\begin{tikzpicture}[scale=1.4]
  % First tetrahedron vertices
  \coordinate (v0) at (-1.5, 0);
  \coordinate (v1) at (0, -0.8);
  \coordinate (v2) at (1.2, 0.8);
  \coordinate (v3) at (0.3, 0.1);

  % Second tetrahedron: new vertex v4
  \coordinate (v4) at (2.8, -0.3);

  % T1 faces
  \fill[dogmalight, opacity=0.3] (v0) -- (v1) -- (v2) -- cycle;
  \fill[dogmalight, opacity=0.2] (v0) -- (v1) -- (v3) -- cycle;

  % Shared face (highlighted)
  \fill[dogmawarm, opacity=0.4] (v1) -- (v2) -- (v3) -- cycle;

  % T2 faces
  \fill[dogmamint, opacity=0.3] (v1) -- (v2) -- (v4) -- cycle;
  \fill[dogmamint, opacity=0.2] (v2) -- (v3) -- (v4) -- cycle;

  % T1 edges
  \draw[thick, dogmablue] (v0) -- (v1);
  \draw[thick, dogmablue] (v0) -- (v2);
  \draw[thick, dogmablue, dashed] (v0) -- (v3);
  \draw[thick, dogmablue] (v1) -- (v2);
  \draw[thick, dogmablue] (v1) -- (v3);
  \draw[thick, dogmablue] (v2) -- (v3);

  % T2 edges (new)
  \draw[very thick, dogmateal] (v1) -- (v4);
  \draw[very thick, dogmateal] (v2) -- (v4);
  \draw[very thick, dogmateal] (v3) -- (v4);

  % Vertices
  \fill[dogmablue] (v0) circle (3pt);
  \fill[dogmablue] (v1) circle (3pt);
  \fill[dogmablue] (v2) circle (3pt);
  \fill[dogmablue] (v3) circle (3pt);
  \fill[dogmateal] (v4) circle (4pt);

  % Labels
  \node[left=4pt, font=\bfseries\small, color=dogmablue] at (v0) {$v_0$};
  \node[below=4pt, font=\bfseries\small, color=dogmablue] at (v1) {$v_1$};
  \node[above=4pt, font=\bfseries\small, color=dogmablue] at (v2) {$v_2$};
  \node[below=4pt, font=\bfseries\small, color=dogmablue] at (v3) {$v_3$};
  \node[right=4pt, font=\bfseries\small, color=dogmateal] at (v4) {$v_4$};

  % Tetrahedron labels
  \node[font=\small\itshape, color=dogmablue] at (-0.5, 0.4) {$\mathcal{T}_1$};
  \node[font=\small\itshape, color=dogmateal] at (2.0, 0.4) {$\mathcal{T}_2$};

  % Shared face annotation
  \draw[thick, dogmared, opacity=0.6, decorate, decoration={brace, amplitude=5pt, mirror}]
    ($(v1)+(0,-0.15)$) -- ($(v3)+(0.1,0.15)$);
  \node[below=8pt, font=\scriptsize, color=dogmared] at (0.4, -0.5) {shared face};
\end{tikzpicture}
\end{center}

\begin{remark}
The shared face $(v_1, v_2, v_3)$ contains three edges and three vertices that
belong to \emph{both} tetrahedra. This is a much richer interface than the
edge-sharing used by triangle meshes: a shared face transmits three pairwise
relationships, while a shared edge transmits only one.
\end{remark}

\subsection{Adding Vertex 5: The Third Tetrahedron}

Continuing the pattern, vertex $v_5$ connects to $v_2, v_3, v_4$, creating:
\[
  \mathcal{T}_3 = (v_2, v_3, v_4, v_5),
\]
which shares face $(v_2, v_3, v_4)$ with $\mathcal{T}_2$.

\subsection{The General Pattern: LinkedTetraChain}

The pattern is now clear. At each step, a new vertex $v_{k+3}$ connects to the
three most recent vertices $(v_k, v_{k+1}, v_{k+2})$, forming tetrahedron:
\[
  \mathcal{T}_{k+1} = (v_k, v_{k+1}, v_{k+2}, v_{k+3}),
\]
which shares face $(v_k, v_{k+1}, v_{k+2})$ with the previous tetrahedron
$\mathcal{T}_k = (v_{k-1}, v_k, v_{k+1}, v_{k+2})$.

\begin{definition}[LinkedTetraChain]
\label{def:linked-tetra-chain}
A \emph{LinkedTetraChain} of length $n$ is an ordered sequence of tetrahedra
$(\mathcal{T}_1, \mathcal{T}_2, \ldots, \mathcal{T}_n)$ constructed from a
vertex sequence $(v_0, v_1, \ldots, v_{n+2})$ where:
\begin{equation}
  \mathcal{T}_k = (v_{k-1}, v_k, v_{k+1}, v_{k+2}), \quad k = 1, 2, \ldots, n,
\end{equation}
and consecutive tetrahedra $\mathcal{T}_k$ and $\mathcal{T}_{k+1}$ share the
triangular face $(v_k, v_{k+1}, v_{k+2})$.
\end{definition}

\begin{proposition}[Edge Count Formula]
\label{prop:edge-count}
A LinkedTetraChain of $n$ tetrahedra built from $n + 3$ vertices has exactly
$3n + 3$ edges.
\end{proposition}
\begin{proof}
The first tetrahedron $\mathcal{T}_1$ contributes $\binom{4}{2} = 6$ edges.
Each subsequent tetrahedron $\mathcal{T}_k$ for $k \geq 2$ adds one new vertex
that connects to 3 existing vertices, contributing exactly 3 new edges.
Thus the total is $6 + 3(n-1) = 3n + 3$.
\end{proof}

The following diagram shows a chain of four linked tetrahedra:

\begin{figure}[htbp]
\centering
\begin{tikzpicture}[scale=1.1, every node/.style={font=\small}]
  % Vertices along the chain
  \coordinate (v0) at (0, 0.2);
  \coordinate (v1) at (1.2, -0.6);
  \coordinate (v2) at (1.0, 0.9);
  \coordinate (v3) at (2.4, 0.2);
  \coordinate (v4) at (3.6, -0.4);
  \coordinate (v5) at (3.8, 1.0);
  \coordinate (v6) at (5.2, 0.1);

  % T1 fill
  \fill[dogmalight, opacity=0.35] (v0) -- (v1) -- (v2) -- cycle;
  \fill[dogmalight, opacity=0.2]  (v0) -- (v1) -- (v3) -- cycle;

  % T2 fill
  \fill[dogmamint, opacity=0.3] (v1) -- (v2) -- (v3) -- cycle;
  \fill[dogmamint, opacity=0.2] (v1) -- (v3) -- (v4) -- cycle;

  % T3 fill
  \fill[dogmawarm, opacity=0.3] (v2) -- (v3) -- (v4) -- cycle;
  \fill[dogmawarm, opacity=0.2] (v3) -- (v4) -- (v5) -- cycle;

  % T4 fill
  \fill[dogmalight, opacity=0.3] (v3) -- (v4) -- (v5) -- cycle;
  \fill[dogmalight, opacity=0.2] (v4) -- (v5) -- (v6) -- cycle;

  % All edges
  \foreach \a/\b in {v0/v1, v0/v2, v0/v3, v1/v2, v1/v3, v2/v3} {
    \draw[thick, dogmablue] (\a) -- (\b);
  }
  \foreach \a/\b in {v1/v4, v2/v4, v3/v4} {
    \draw[thick, dogmateal] (\a) -- (\b);
  }
  \foreach \a/\b in {v2/v5, v3/v5, v4/v5} {
    \draw[thick, dogmagold] (\a) -- (\b);
  }
  \foreach \a/\b in {v3/v6, v4/v6, v5/v6} {
    \draw[thick, dogmared] (\a) -- (\b);
  }

  % Vertices
  \foreach \v in {v0, v1, v2, v3, v4, v5, v6} {
    \fill[dogmablue!80!black] (\v) circle (3pt);
  }

  % Vertex labels
  \node[left=3pt] at (v0) {$v_0$};
  \node[below=3pt] at (v1) {$v_1$};
  \node[above=3pt] at (v2) {$v_2$};
  \node[below=3pt] at (v3) {$v_3$};
  \node[below=3pt] at (v4) {$v_4$};
  \node[above=3pt] at (v5) {$v_5$};
  \node[right=3pt] at (v6) {$v_6$};

  % Tetrahedron labels (below)
  \node[font=\small\bfseries, color=dogmablue] at (0.8, -1.1) {$\mathcal{T}_1$};
  \node[font=\small\bfseries, color=dogmateal] at (2.0, -1.1) {$\mathcal{T}_2$};
  \node[font=\small\bfseries, color=dogmagold] at (3.4, -1.1) {$\mathcal{T}_3$};
  \node[font=\small\bfseries, color=dogmared] at (4.8, -1.1) {$\mathcal{T}_4$};
\end{tikzpicture}
\caption{A LinkedTetraChain of four tetrahedra. Each tetrahedron shares a
triangular face with its neighbors. Colors distinguish the edges added by each
new vertex. Total: 7 vertices, $3(4)+3 = 15$ edges.}
\label{fig:linked-chain}
\end{figure}

\begin{keyinsight}[The Sliding Window Analogy]
The LinkedTetraChain is equivalent to a \emph{sliding window of width 4 with
overlap 3} over the vertex sequence. Each tetrahedron ``sees'' four consecutive
vertices, and consecutive tetrahedra overlap in three vertices. This means that
each vertex participates in up to four tetrahedra (except near the ends),
ensuring that information about each vertex is propagated through multiple
local interaction contexts.

Compare this to the sliding window in signal processing: a window of width $W$
with overlap $W - 1$ slides one position at a time. Here, $W = 4$ and
overlap $= 3$, so each step advances by one vertex. But unlike a flat sliding
window, each ``window'' is a tetrahedron with volume, chirality, and full $K_4$
connectivity.
\end{keyinsight}

\subsection{Growth Table}

The following table tracks the combinatorics as we add vertices:

\begin{center}
\renewcommand{\arraystretch}{1.15}
\begin{tabular}{ccccccl}
\toprule
\textbf{Vertex} & \textbf{Total $|V|$} & \textbf{New edges} & \textbf{Total $|E|$}
  & \textbf{Total $|\mathcal{T}|$} & \textbf{Shared face} & \textbf{Graph} \\
\midrule
$v_0$         & 1 & 0 & 0  & 0 & --- & $K_1$ \\
$v_1$         & 2 & 1 & 1  & 0 & --- & $K_2$ \\
$v_2$         & 3 & 2 & 3  & 0 & --- & $K_3$ \\
$v_3$         & 4 & 3 & 6  & 1 ($\mathcal{T}_1$) & --- & $K_4$ \\
$v_4$         & 5 & 3 & 9  & 2 ($\mathcal{T}_2$) & $(v_1,v_2,v_3)$ & --- \\
$v_5$         & 6 & 3 & 12 & 3 ($\mathcal{T}_3$) & $(v_2,v_3,v_4)$ & --- \\
$v_6$         & 7 & 3 & 15 & 4 ($\mathcal{T}_4$) & $(v_3,v_4,v_5)$ & --- \\
$v_{k+3}$     & $k{+}4$ & 3 & $3(k{+}1){+}3$ & $k{+}1$ & $(v_k,v_{k+1},v_{k+2})$ & --- \\
\bottomrule
\end{tabular}
\end{center}

Notice that after the first tetrahedron is formed, each new vertex adds exactly
3 new edges and exactly 1 new tetrahedron. The growth is perfectly linear.

% ═══════════════════════════════════════════════════════════════════════════════
\section{DNA Mapping --- The Key Insight}
\label{sec:dna-mapping}
% ═══════════════════════════════════════════════════════════════════════════════

We now arrive at the central insight of the TetraData Structure: the mapping
between tetrahedra and DNA sequences.

\subsection{Four Vertices, Four Bases}

A tetrahedron has four vertices. DNA has four bases: adenine (A), cytosine (C),
guanine (G), and thymine (T). In a LinkedTetraChain, each vertex is assigned
one base type from $\{\mathsf{A}, \mathsf{C}, \mathsf{G}, \mathsf{T}\}$.

\begin{dogmadef}[TetraData--DNA Correspondence]
In the TetraData Structure, each vertex $v_i$ carries a base assignment
$b_i \in \{\mathsf{A}, \mathsf{C}, \mathsf{G}, \mathsf{T}\}$.
A tetrahedron $\mathcal{T}_k = (v_{k-1}, v_k, v_{k+1}, v_{k+2})$ with base
assignments $(b_{k-1}, b_k, b_{k+1}, b_{k+2})$ represents a \emph{window} of
four consecutive bases in a DNA-like sequence.
\end{dogmadef}

\subsection{Worked Example: Mapping \texorpdfstring{$5'$}{5'}-ACGTACGT-\texorpdfstring{$3'$}{3'}}

Consider the DNA sequence $5'$-\texttt{ACGTACGT}-$3'$, which has 8 bases. We
assign these to vertices $v_0$ through $v_7$:

\begin{center}
\renewcommand{\arraystretch}{1.15}
\begin{tabular}{ccccccccc}
\toprule
Vertex & $v_0$ & $v_1$ & $v_2$ & $v_3$ & $v_4$ & $v_5$ & $v_6$ & $v_7$ \\
\midrule
Base   & \textsf{A} & \textsf{C} & \textsf{G} & \textsf{T} & \textsf{A} & \textsf{C} & \textsf{G} & \textsf{T} \\
\bottomrule
\end{tabular}
\end{center}

This sequence of 8 vertices produces a LinkedTetraChain of $8 - 3 = 5$ tetrahedra:

\begin{center}
\renewcommand{\arraystretch}{1.15}
\begin{tabular}{clcl}
\toprule
\textbf{Tetra} & \textbf{Vertices} & \textbf{Bases} & \textbf{Shared face with prev.} \\
\midrule
$\mathcal{T}_1$ & $(v_0, v_1, v_2, v_3)$ & (A, C, G, T) & --- \\
$\mathcal{T}_2$ & $(v_1, v_2, v_3, v_4)$ & (C, G, T, A) & $(v_1, v_2, v_3) = $ (C, G, T) \\
$\mathcal{T}_3$ & $(v_2, v_3, v_4, v_5)$ & (G, T, A, C) & $(v_2, v_3, v_4) = $ (G, T, A) \\
$\mathcal{T}_4$ & $(v_3, v_4, v_5, v_6)$ & (T, A, C, G) & $(v_3, v_4, v_5) = $ (T, A, C) \\
$\mathcal{T}_5$ & $(v_4, v_5, v_6, v_7)$ & (A, C, G, T) & $(v_4, v_5, v_6) = $ (A, C, G) \\
\bottomrule
\end{tabular}
\end{center}

\begin{figure}[htbp]
\centering
\begin{tikzpicture}[scale=1.0, every node/.style={font=\small}]
  % DNA sequence on top
  \foreach \i/\base/\col in {0/A/dogmared, 1/C/dogmablue, 2/G/dogmateal, 3/T/dogmagold,
                              4/A/dogmared, 5/C/dogmablue, 6/G/dogmateal, 7/T/dogmagold} {
    \node[circle, draw=\col, fill=\col!15, minimum size=20pt, font=\bfseries\small]
      (b\i) at (1.3*\i, 2.5) {\base};
    \node[below=1pt, font=\tiny, color=black!60] at (b\i.south) {$v_\i$};
  }

  % Edges connecting consecutive bases
  \foreach \i [evaluate=\i as \j using int(\i+1)] in {0,...,6} {
    \draw[thick, black!30] (b\i) -- (b\j);
  }

  % Tetrahedra brackets below
  \foreach \k/\s/\col in {1/0/dogmared, 2/1/dogmablue, 3/2/dogmateal, 4/3/dogmagold, 5/4/dogmared} {
    \pgfmathsetmacro{\left}{1.3*\s - 0.3}
    \pgfmathsetmacro{\right}{1.3*(\s+3) + 0.3}
    \pgfmathsetmacro{\yy}{1.2 - 0.35*(\k-1)}
    \draw[thick, \col, decorate, decoration={brace, amplitude=4pt, mirror}]
      (\left, \yy) -- (\right, \yy)
      node[midway, below=5pt, font=\scriptsize\bfseries, color=\col]
      {$\mathcal{T}_\k$};
  }
\end{tikzpicture}
\caption{DNA sequence \texttt{ACGTACGT} mapped to a LinkedTetraChain.
Each brace shows which four consecutive vertices form a tetrahedron.
Note the overlap: each tetrahedron shares three vertices with its neighbors.}
\label{fig:dna-mapping}
\end{figure}

Notice the structure: each tetrahedron is a sliding window of width 4 that
advances by one base at a time. The overlap of 3 bases ensures that
\emph{every pair of adjacent bases} is captured within at least one tetrahedron,
and in fact every pair of bases separated by at most 2 positions is captured.

\subsection{Why Not Triangles? An Information Density Argument}

If we used triangles instead of tetrahedra, our sliding window would be width 3
with overlap 2. Each triangle would be a $K_3$ with only 3 edges. Let us
compare the two approaches for the same 8-base sequence:

\begin{center}
\renewcommand{\arraystretch}{1.15}
\begin{tabular}{lcc}
\toprule
& \textbf{Triangle mesh} ($K_3$) & \textbf{Tetra mesh} ($K_4$) \\
\midrule
Window width & 3 & 4 \\
Overlap & 2 & 3 \\
Elements for 8 bases & 6 triangles & 5 tetrahedra \\
Edges per element & 3 & 6 \\
Total unique edges & 7 & 10 \\
Pairwise interactions & 3 per element & 6 per element \\
Has volume? & No & Yes \\
Has chirality? & No & Yes \\
\bottomrule
\end{tabular}
\end{center}

The tetrahedron captures $2\times$ the pairwise interactions per element, while
also encoding volume and chirality. For the same sequence, the tetrahedral
representation is strictly richer.

\subsection{The Six Edges: Pairwise Base Interactions}
\label{sec:six-edges}

Within each tetrahedron $\mathcal{T}_k = (b_{k-1}, b_k, b_{k+1}, b_{k+2})$, the six edges
encode all pairwise interactions among the four bases:

\begin{center}
\begin{tikzpicture}[scale=1.5]
  % Tetrahedron vertices with base labels
  \coordinate (A) at (0,0);
  \coordinate (C) at (2.5,0);
  \coordinate (G) at (1.25,2.0);
  \coordinate (T) at (0.8,0.7);

  % Faces
  \fill[dogmalight, opacity=0.3] (A) -- (C) -- (G) -- cycle;
  \fill[dogmamint, opacity=0.3] (C) -- (G) -- (T) -- cycle;

  % Watson-Crick edges (thick, colored)
  \draw[very thick, dogmared] (A) -- (T)
    node[midway, left=3pt, font=\small\bfseries, color=dogmared] {A--T};
  \draw[very thick, dogmateal] (C) -- (G)
    node[midway, right=3pt, font=\small\bfseries, color=dogmateal] {C--G};

  % Other edges
  \draw[thick, dogmablue!60] (A) -- (C)
    node[midway, below=2pt, font=\scriptsize] {A--C};
  \draw[thick, dogmablue!60] (A) -- (G)
    node[midway, above left=1pt, font=\scriptsize] {A--G};
  \draw[thick, dogmablue!60] (C) -- (T)
    node[midway, below right=1pt, font=\scriptsize] {C--T};
  \draw[thick, dogmablue!60] (G) -- (T)
    node[midway, above right=1pt, font=\scriptsize] {G--T};

  % Vertices
  \fill[dogmared] (A) circle (4pt);
  \fill[dogmablue] (C) circle (4pt);
  \fill[dogmateal] (G) circle (4pt);
  \fill[dogmagold] (T) circle (4pt);

  % Labels
  \node[below left=3pt, font=\bfseries, color=dogmared] at (A) {A};
  \node[below right=3pt, font=\bfseries, color=dogmablue] at (C) {C};
  \node[above=4pt, font=\bfseries, color=dogmateal] at (G) {G};
  \node[below=4pt, font=\bfseries, color=dogmagold] at (T) {T};

  % Legend
  \node[right=20pt, align=left, font=\small] at (2.5, 1.5) {
    \textcolor{dogmared}{\rule{12pt}{2pt}} Watson--Crick (A--T) \\[2pt]
    \textcolor{dogmateal}{\rule{12pt}{2pt}} Watson--Crick (C--G) \\[2pt]
    \textcolor{dogmablue!60}{\rule{12pt}{2pt}} Other pairwise
  };
\end{tikzpicture}
\end{center}

The six edges naturally partition into:
\begin{itemize}[leftmargin=1.8em]
  \item \textbf{2 Watson--Crick edges:} A--T and C--G. These are the
        complementary base pairings that form hydrogen bonds in the double helix.
  \item \textbf{4 non-Watson--Crick edges:} A--C, A--G, C--T, G--T. These
        encode stacking interactions, non-canonical base pairing, and other
        local sequence features.
\end{itemize}

Each edge carries a learnable weight $w_{ij} \in \mathbb{R}$ that represents
the strength and character of the pairwise interaction. These weights are
analogous to attention scores in a transformer, but they are \emph{local}
(within a single tetrahedron) and \emph{structured} (constrained by the
tetrahedral geometry).

% ═══════════════════════════════════════════════════════════════════════════════
\section{The TetraData Structure (TDS) --- Formal Definition}
\label{sec:tds}
% ═══════════════════════════════════════════════════════════════════════════════

We now give the complete formal definition of the TetraData Structure,
incorporating everything we have built up.

\begin{definition}[TetraData Structure]
\label{def:tds}
A \emph{TetraData Structure} (TDS) element is defined by:
\begin{equation}
\mathcal{T}_i = (\mathbf{V}_i, \mathbf{E}_i, \mathbf{B}_i, v_i, \chi_i),
\end{equation}
where:
\begin{itemize}[leftmargin=1.5em]
  \item $\mathbf{V}_i \in \mathbb{R}^{4 \times 3}$: four vertex positions in 3D space.
  \item $\mathbf{E}_i \in \mathbb{R}^{6}$: six edge weights (forming the complete graph $K_4$).
  \item $\mathbf{B}_i \in \{\mathsf{A}, \mathsf{C}, \mathsf{G}, \mathsf{T}\}^4$: base type assignments for each vertex.
  \item $v_i \in \mathbb{R}$: enclosed volume (computed from vertex positions).
  \item $\chi_i \in \{+1, -1\}$: chirality (handedness), defined by the sign of the determinant.
\end{itemize}
\end{definition}

\subsection{Chirality}
\label{sec:chirality}

\emph{Chirality}---the property of non-superimposability on a mirror image---is
fundamental to molecular biology. Amino acids are left-handed ($L$-amino acids),
sugars are right-handed ($D$-sugars), and the DNA double helix is right-handed.
TDS preserves chirality through the sign of the vertex determinant:

\begin{equation}
\chi_i = \text{sign}\left(\det \begin{pmatrix} \mathbf{v}_1 - \mathbf{v}_0 & \mathbf{v}_2 - \mathbf{v}_0 & \mathbf{v}_3 - \mathbf{v}_0 \end{pmatrix}\right).
\end{equation}

Chirality-preserving transformations are those that maintain the sign of $\chi$.
This constraint ensures that biological structures retain their handedness
through computational operations.

\begin{example}[Computing Chirality]
\label{ex:chirality}
Consider a tetrahedron with vertices:
\[
  \mathbf{v}_0 = (0,0,0), \quad
  \mathbf{v}_1 = (1,0,0), \quad
  \mathbf{v}_2 = (0,1,0), \quad
  \mathbf{v}_3 = (0,0,1).
\]
The edge matrix is:
\[
  M = \begin{pmatrix}
    1-0 & 0-0 & 0-0 \\
    0-0 & 1-0 & 0-0 \\
    0-0 & 0-0 & 1-0
  \end{pmatrix}
  = \begin{pmatrix} 1 & 0 & 0 \\ 0 & 1 & 0 \\ 0 & 0 & 1 \end{pmatrix} = I_3.
\]
We have $\det(M) = +1$, so $\chi = +1$ (right-handed). The signed volume is
$V_{\text{signed}} = \frac{1}{6}(+1) = +\frac{1}{6}$.

Now reflect $\mathbf{v}_3$ across the $xy$-plane: $\mathbf{v}_3' = (0,0,-1)$.
The new determinant is $-1$, giving $\chi = -1$ (left-handed). Reflection
reverses chirality, as expected.
\end{example}

\subsection{Face-Sharing Adjacency}

Adjacent tetrahedra share triangular faces, forming \emph{tetrahedral chains}
or \emph{tetrahedral meshes}. This face-sharing adjacency defines a natural
graph structure:

\begin{definition}[Tetrahedral Mesh]
A \emph{tetrahedral mesh} is a collection of tetrahedra
$\{\mathcal{T}_1, \ldots, \mathcal{T}_M\}$ where adjacent elements share
triangular faces. The \emph{dual graph} has one node per tetrahedron and edges
between face-sharing neighbors.
\end{definition}

Face-sharing adjacency is richer than the edge-sharing adjacency of triangle
meshes: each tetrahedron has exactly four potential neighbors (one per face),
creating a bounded-degree graph that enables efficient local computation.

% ═══════════════════════════════════════════════════════════════════════════════
\section{36 Degrees of Freedom}
\label{sec:36-dof}
% ═══════════════════════════════════════════════════════════════════════════════

Each TDS element carries a rich set of information. Let us count the degrees
of freedom precisely.

\subsection{Detailed DOF Breakdown}

\begin{center}
\renewcommand{\arraystretch}{1.25}
\begin{tabular}{clcl}
\toprule
\textbf{DOF} & \textbf{Component} & \textbf{Count} & \textbf{Description} \\
\midrule
1--12 & Vertex positions $\mathbf{V}_i \in \mathbb{R}^{4 \times 3}$ & 12
  & 4 vertices $\times$ 3 coordinates each \\
13--18 & Edge weights $\mathbf{E}_i \in \mathbb{R}^6$ & 6
  & All $\binom{4}{2} = 6$ pairwise interactions \\
19--34 & Base encoding $\text{onehot}(\mathbf{B}_i) \in \mathbb{R}^{4 \times 4}$ & 16
  & 4 vertices $\times$ 4-dim one-hot encoding \\
35 & Volume $v_i \in \mathbb{R}$ & 1
  & Signed volume of the tetrahedron \\
36 & Chirality $\chi_i \in \{+1,-1\}$ & 1
  & Handedness (sign of the determinant) \\
\midrule
& \textbf{Total} & \textbf{36} & \\
\bottomrule
\end{tabular}
\end{center}

\subsection{Comparison: Triangle vs.\ Tetrahedron DOF}

\begin{center}
\renewcommand{\arraystretch}{1.15}
\begin{tabular}{lcc}
\toprule
\textbf{Component} & \textbf{Triangle} & \textbf{Tetrahedron} \\
\midrule
Vertex positions & $3 \times 3 = 9$ & $4 \times 3 = 12$ \\
Edge weights & $\binom{3}{2} = 3$ & $\binom{4}{2} = 6$ \\
Base encoding & --- & $4 \times 4 = 16$ \\
Volume & 0 (area only) & 1 \\
Chirality & 0 (no handedness) & 1 \\
\midrule
\textbf{Total DOF} & \textbf{12} & \textbf{36} \\
\midrule
\textbf{Ratio} & \multicolumn{2}{c}{$36/12 = 3\times$ richer} \\
\bottomrule
\end{tabular}
\end{center}

If we count only the positional DOF (no base encoding), the triangle has 9 DOF
and the tetrahedron has $12 + 6 + 1 + 1 = 20$ DOF---still more than $2\times$
richer. The base encoding adds another 16 DOF that are entirely absent in
triangle meshes.

\subsection{TetraTokenization}
\label{sec:tetra-tokenization}

For foundation model integration, each TDS element is serialized into a
36-dimensional vector:

\begin{equation}
\text{TetraToken}(\mathcal{T}_i) = [\underbrace{\mathbf{V}_i^{\text{flat}}}_{12} \;|\; \underbrace{\mathbf{E}_i}_{6} \;|\; \underbrace{\text{onehot}(\mathbf{B}_i)}_{16} \;|\; \underbrace{v_i}_{1} \;|\; \underbrace{\chi_i}_{1}] \in \mathbb{R}^{36}.
\label{eq:tetra-token}
\end{equation}

This tokenization enables DOGMA foundation models to ingest and produce 3D
structural data alongside text, code, and genomic sequences.

\begin{implnote}[TetraToken in Practice]
In the DOGMA-Supreme implementation, the 36-dimensional TetraToken is projected
to the model's hidden dimension $d_{\text{model}}$ via a learned linear layer:
\[
  \mathbf{h}_i = \mathbf{W}_{\text{proj}} \cdot \text{TetraToken}(\mathcal{T}_i) + \mathbf{b}_{\text{proj}},
  \quad \mathbf{W}_{\text{proj}} \in \mathbb{R}^{d_{\text{model}} \times 36}.
\]
The reverse direction (decoding) uses a symmetric linear projection from
$d_{\text{model}}$ back to 36 dimensions, followed by extraction of vertex
positions, edge weights, base assignments, volume, and chirality.
\end{implnote}

% ═══════════════════════════════════════════════════════════════════════════════
\section{TetraMemory: Structured Memory for DOGMA}
\label{sec:tetradata-tetramemory}
% ═══════════════════════════════════════════════════════════════════════════════

TetraMemory is DOGMA's alternative to self-attention. Instead of computing
pairwise interactions between all positions (quadratic cost), TetraMemory
organizes state elements into a tetrahedral mesh where interactions are
\emph{local but compose into global computation} through propagation.

\subsection{The Quadratic Attention Problem}
\label{sec:quadratic-problem}

Standard self-attention computes:
\begin{equation}
\text{Attention}(\mathbf{Q}, \mathbf{K}, \mathbf{V}) = \text{softmax}\left(\frac{\mathbf{Q}\mathbf{K}^\top}{\sqrt{d_k}}\right)\mathbf{V}, \quad \text{cost: } O(T^2 d).
\end{equation}

For a sequence of length $T = 100{,}000$ (a modest genomic sequence), the
attention matrix has $10^{10}$ entries---impractical for both memory and
computation. Even for $T = 10{,}000$, storing the $T \times T$ attention matrix
in float32 requires $\sim$400 MB per head, per layer.

\begin{center}
\begin{tikzpicture}[scale=0.8]
  % Self-attention: all-to-all connections
  \begin{scope}[shift={(0,0)}]
    \node[font=\bfseries\small, color=dogmablue] at (1.5, 3.2) {Self-Attention};
    \node[font=\scriptsize, color=dogmagray] at (1.5, 2.7) {$O(T^2)$ connections};

    \foreach \i in {0,...,5} {
      \fill[dogmablue] (0.5*\i, 0) circle (3pt);
      \fill[dogmablue] (0.5*\i, 2) circle (3pt);
    }
    % All-to-all connections (subset shown)
    \foreach \i in {0,...,5} {
      \foreach \j in {0,...,5} {
        \draw[dogmablue!20, thin] (0.5*\i, 0) -- (0.5*\j, 2);
      }
    }
  \end{scope}

  % TetraMemory: local connections
  \begin{scope}[shift={(6,0)}]
    \node[font=\bfseries\small, color=dogmateal] at (1.5, 3.2) {TetraMemory};
    \node[font=\scriptsize, color=dogmagray] at (1.5, 2.7) {$O(T)$ connections};

    \foreach \i in {0,...,5} {
      \fill[dogmateal] (0.5*\i, 0) circle (3pt);
      \fill[dogmateal] (0.5*\i, 2) circle (3pt);
    }
    % Local connections only (window of 4)
    \foreach \i in {0,...,5} {
      \pgfmathsetmacro{\lo}{max(0, \i-1)}
      \pgfmathsetmacro{\hi}{min(5, \i+2)}
      \foreach \j in {\lo,...,\hi} {
        \draw[dogmateal!50, thin] (0.5*\i, 0) -- (0.5*\j, 2);
      }
    }
  \end{scope}
\end{tikzpicture}
\end{center}

\subsection{Face Propagation: Information Flow Through Shared Faces}
\label{sec:face-propagation}

The key mechanism of TetraMemory is \emph{face propagation}: information
flows from one tetrahedron to its neighbor through their shared triangular face.

In a LinkedTetraChain, $\mathcal{T}_k$ and $\mathcal{T}_{k+1}$ share three
vertices. When we compute the state of $\mathcal{T}_{k+1}$, three of its four
vertices already carry information from $\mathcal{T}_k$. Only one vertex is
new. This means that 75\% of each tetrahedron's information comes from its
predecessor---a natural form of \emph{causal propagation}.

\begin{definition}[TetraMemory State Update]
\label{def:tetra-update}
At layer $\ell$, the state of tetrahedron $\mathcal{T}_i$ is updated by:
\begin{equation}
\mathbf{h}_i^{(\ell+1)} = \phi\left(\mathbf{h}_i^{(\ell)} + \sum_{j \in \mathcal{N}_4(i)} \alpha_{ij} \cdot \mathbf{W}^{(\ell)} \mathbf{h}_j^{(\ell)}\right),
\label{eq:tetra-update}
\end{equation}
where:
\begin{itemize}[leftmargin=1.5em]
  \item $\mathcal{N}_4(i)$ is the set of (at most 4) face-sharing neighbors of $\mathcal{T}_i$.
  \item $\alpha_{ij}$ are attention weights computed from edge properties and shared-face geometry.
  \item $\mathbf{W}^{(\ell)} \in \mathbb{R}^{d \times d}$ is the layer's learnable weight matrix.
  \item $\phi$ is a nonlinear activation function.
\end{itemize}
\end{definition}

The cost per layer is $O(T \cdot k \cdot d)$ where $k \leq 4$, making it
\emph{linear} in sequence length.

\subsection{Volume as Attention Weight}
\label{sec:volume-attention}

Each tetrahedron's signed volume provides a natural attention modulator:

\begin{equation}
\alpha_{ij} = \sigma\!\left(\frac{\mathbf{q}_i^\top \mathbf{k}_j}{\sqrt{d}} + \beta \cdot V_{\text{signed}}(\mathcal{T}_i)\right),
\end{equation}

where $\sigma$ is a sigmoid function and $\beta$ is a learnable temperature
parameter. The volume term acts as a \emph{geometric bias}: tetrahedra with
larger volume (more ``open'' configurations) attend more strongly to their
neighbors, while degenerate (near-flat) tetrahedra attend weakly.

\begin{bioinsight}[Volume as Importance]
In molecular biology, the volume of a protein's active site determines its
binding affinity. A large, open active site can accommodate many substrates;
a narrow, constrained one is highly specific. TetraMemory's use of volume as
an attention modulator mirrors this: the ``openness'' of a tetrahedral
configuration controls how much information flows through it.
\end{bioinsight}

\subsection{Edge Weights as Pairwise Compatibility}
\label{sec:edge-compatibility}

The six edge weights $\mathbf{E}_i \in \mathbb{R}^6$ within each tetrahedron
encode pairwise compatibility scores between the four vertices. These are
analogous to the pairwise attention scores in self-attention, but:

\begin{enumerate}[leftmargin=2em]
  \item They are \emph{local}: only the 6 pairs within a tetrahedron are
        computed, not the $\binom{T}{2}$ global pairs.
  \item They are \emph{geometrically constrained}: the edge weights must be
        consistent with a valid 3D tetrahedron (triangle inequality on faces,
        positive volume).
  \item They are \emph{typed}: the Watson--Crick edges (A--T, C--G) can carry
        different learned biases than the non-Watson--Crick edges.
\end{enumerate}

\subsection{Long-Range Information Flow}
\label{sec:long-range}

How does local interaction enable global computation? Through \emph{propagation}:
after $L$ layers of TetraMemory updates, information from tetrahedron
$\mathcal{T}_i$ has propagated to all tetrahedra within graph distance $L$ in
the dual graph.

For a LinkedTetraChain of $n$ tetrahedra, the dual graph is a path graph of
length $n$, so the diameter is $n$ and full propagation requires $L = n$ layers.
However, for a 3D tetrahedral mesh (not just a chain), the dual graph diameter
is $O(n^{1/3})$, so $O(n^{1/3})$ layers suffice for global information flow.

\begin{theorem}[TetraMemory Propagation]
\label{thm:propagation}
For a tetrahedral mesh with $n$ tetrahedra and dual-graph diameter $D$, after
$L \geq D$ layers of TetraMemory updates (Equation~\ref{eq:tetra-update}),
every tetrahedron has received information from every other tetrahedron.
\end{theorem}
\begin{proof}
At each layer, each tetrahedron receives information from its face-sharing
neighbors (at most 4). After $L$ layers, the receptive field of each
tetrahedron is the set of all tetrahedra within graph distance $L$ in the
dual graph. When $L \geq D$, this includes all tetrahedra.
\end{proof}

\subsection{Complexity Comparison}

\begin{center}
\renewcommand{\arraystretch}{1.25}
\begin{tabular}{lcc}
\toprule
& \textbf{Self-Attention} & \textbf{TetraMemory} \\
\midrule
Cost per layer & $O(T^2 d)$ & $O(T \cdot k \cdot d)$, $k \leq 4$ \\
Memory per layer & $O(T^2)$ & $O(T \cdot k)$ \\
Layers for global info & 1 (instant) & $O(D)$ where $D$ is mesh diameter \\
Total cost for global & $O(T^2 d)$ & $O(D \cdot T \cdot k \cdot d)$ \\
For 3D mesh ($D = O(T^{1/3})$) & $O(T^2 d)$ & $O(T^{4/3} k d)$ \\
Structural locality & None & Built-in (mesh topology) \\
Chirality encoding & None & Built-in (signed volume) \\
\bottomrule
\end{tabular}
\end{center}

\begin{keyinsight}[TetraMemory vs.\ Attention: The Biological Argument]
Biological neural networks do not use all-to-all connectivity. Instead, neurons
interact with local neighbors, and global computation emerges from multi-hop
signal propagation through network topology. TetraMemory follows this principle:
local interactions compose into global computation, just as local gene
regulatory interactions compose into genome-wide expression patterns.

The key advantage is not just computational efficiency. It is \emph{structural
locality}: the mesh topology encodes which elements should interact most
strongly, providing an architectural prior that dense attention lacks.
\end{keyinsight}

% ═══════════════════════════════════════════════════════════════════════════════
\section{The DoubleHelixMesh}
\label{sec:double-helix-mesh}
% ═══════════════════════════════════════════════════════════════════════════════

In real DNA, information is stored in two complementary strands running
antiparallel to each other, forming the iconic double helix. DOGMA mirrors this
with the \textbf{DoubleHelixMesh}: two complementary LinkedTetraChains---the
\emph{sense} strand and the \emph{antisense} strand---cross-linked by
Watson--Crick hydrogen bonds.

\subsection{Sense and Antisense Strands}

Given a sense-strand sequence $S = s_0 s_1 s_2 \cdots s_{n-1}$, the antisense
strand is the reverse complement:
\begin{equation}
  \bar{S} = \overline{s_{n-1}}\;\overline{s_{n-2}}\;\cdots\;\overline{s_1}\;\overline{s_0},
\end{equation}
where the complement map is $\overline{\mathsf{A}} = \mathsf{T}$,
$\overline{\mathsf{T}} = \mathsf{A}$, $\overline{\mathsf{C}} = \mathsf{G}$,
$\overline{\mathsf{G}} = \mathsf{C}$.

Each strand forms its own LinkedTetraChain:
\begin{itemize}[leftmargin=1.8em]
  \item \textbf{Sense chain:} $\mathcal{C}_{\text{sense}} = (\mathcal{T}_1^{+}, \mathcal{T}_2^{+}, \ldots, \mathcal{T}_{n-3}^{+})$
  \item \textbf{Antisense chain:} $\mathcal{C}_{\text{anti}} = (\mathcal{T}_1^{-}, \mathcal{T}_2^{-}, \ldots, \mathcal{T}_{n-3}^{-})$
\end{itemize}

The antisense chain runs in the reverse direction (antiparallel), so
$\mathcal{T}_k^{-}$ aligns spatially with $\mathcal{T}_{n-3-k}^{+}$.

\subsection{Watson--Crick Cross-Links}

The two strands are connected by hydrogen bonds between complementary base pairs:

\begin{itemize}[leftmargin=1.8em]
  \item \textbf{A--T pairing:} 2 hydrogen bonds (weaker).
  \item \textbf{G--C pairing:} 3 hydrogen bonds (stronger).
\end{itemize}

In the DoubleHelixMesh, these cross-links are represented as inter-chain edges
with weights proportional to hydrogen bond strength:

\begin{equation}
  w_{\text{cross}}(b, \bar{b}) = \begin{cases}
    w_{\text{AT}} & \text{if } (b, \bar{b}) \in \{(\mathsf{A}, \mathsf{T}), (\mathsf{T}, \mathsf{A})\}, \\
    w_{\text{GC}} & \text{if } (b, \bar{b}) \in \{(\mathsf{G}, \mathsf{C}), (\mathsf{C}, \mathsf{G})\},
  \end{cases}
\end{equation}

where $w_{\text{GC}} > w_{\text{AT}}$ reflects the stronger G--C bond.

\begin{figure}[htbp]
\centering
\begin{tikzpicture}[scale=0.95, every node/.style={font=\small}]
  % Sense strand (top)
  \node[font=\bfseries\small, color=dogmablue] at (-1.5, 1.8) {Sense $5' \to 3'$};
  \foreach \i/\base/\col in {0/A/dogmared, 1/C/dogmablue, 2/G/dogmateal,
                              3/T/dogmagold, 4/A/dogmared, 5/C/dogmablue,
                              6/G/dogmateal, 7/T/dogmagold} {
    \node[circle, draw=\col, fill=\col!15, minimum size=18pt,
          inner sep=0pt, font=\bfseries\scriptsize]
      (s\i) at (1.2*\i, 1.5) {\base};
  }
  % Sense strand backbone
  \foreach \i [evaluate=\i as \j using int(\i+1)] in {0,...,6} {
    \draw[thick, dogmablue!40] (s\i) -- (s\j);
  }

  % Antisense strand (bottom, reversed complement)
  \node[font=\bfseries\small, color=dogmateal] at (-1.5, -1.8) {Antisense $3' \to 5'$};
  \foreach \i/\base/\col in {0/T/dogmagold, 1/G/dogmateal, 2/C/dogmablue,
                              3/A/dogmared, 4/T/dogmagold, 5/G/dogmateal,
                              6/C/dogmablue, 7/A/dogmared} {
    \node[circle, draw=\col, fill=\col!15, minimum size=18pt,
          inner sep=0pt, font=\bfseries\scriptsize]
      (a\i) at (1.2*\i, -1.5) {\base};
  }
  % Antisense backbone
  \foreach \i [evaluate=\i as \j using int(\i+1)] in {0,...,6} {
    \draw[thick, dogmateal!40] (a\i) -- (a\j);
  }

  % Watson-Crick hydrogen bonds
  \foreach \i/\nbonds in {0/2, 1/3, 2/3, 3/2, 4/2, 5/3, 6/3, 7/2} {
    \pgfmathsetmacro{\aindex}{7 - \i}
    \ifnum\nbonds=2
      % A-T: 2 bonds (dashed)
      \draw[dogmared!60, thick, densely dashed] (s\i) -- (a\aindex);
      \node[font=\tiny, color=dogmared!60, right=1pt] at ($(s\i)!0.5!(a\aindex)$) {2H};
    \else
      % G-C: 3 bonds (solid)
      \draw[dogmateal!60, thick] (s\i) -- (a\aindex);
      \node[font=\tiny, color=dogmateal!60, right=1pt] at ($(s\i)!0.5!(a\aindex)$) {3H};
    \fi
  }
\end{tikzpicture}
\caption{The DoubleHelixMesh for sequence \texttt{ACGTACGT}. The sense strand
(top, $5' \to 3'$) and antisense strand (bottom, $3' \to 5'$) are connected by
Watson--Crick hydrogen bonds. A--T pairs have 2 hydrogen bonds (dashed);
G--C pairs have 3 hydrogen bonds (solid).}
\label{fig:double-helix}
\end{figure}

\subsection{Advantages of the Double Helix Structure}

The DoubleHelixMesh provides several computational advantages:

\begin{enumerate}[leftmargin=2em]
  \item \textbf{Redundancy and error detection:} Each base on the sense strand
        has a complementary base on the antisense strand. Discrepancies can be
        detected and corrected, mirroring DNA's own error-correction mechanism.

  \item \textbf{Bidirectional reading:} The sense strand reads $5' \to 3'$ while
        the antisense reads $3' \to 5'$. This provides a natural mechanism for
        bidirectional sequence processing without the need for separate forward
        and backward passes.

  \item \textbf{Cross-strand attention:} The hydrogen bond cross-links create
        additional attention pathways between the two chains, enriching the
        information flow beyond what a single chain provides.

  \item \textbf{Melting dynamics:} In biology, the double helix can ``melt''
        (denature) at high temperatures, separating the strands. In DOGMA, a
        learned ``temperature'' parameter controls the strength of cross-strand
        interactions, enabling the model to dynamically decouple the strands
        when independent processing is beneficial.
\end{enumerate}

\begin{implnote}[DoubleHelixMesh Efficiency]
Despite having two chains, the DoubleHelixMesh does \emph{not} double the
computational cost. The antisense chain shares vertex positions (by symmetry)
with the sense chain, and the cross-strand hydrogen bond weights are computed
from the base assignments (which are determined by complementarity). The
additional cost is $O(T)$ for the cross-strand edges, which is dominated by
the $O(T \cdot k \cdot d)$ per-layer cost of TetraMemory updates on each chain.
\end{implnote}

% ═══════════════════════════════════════════════════════════════════════════════
\section{Delaunay Tetrahedralization}
\label{sec:delaunay}
% ═══════════════════════════════════════════════════════════════════════════════

While the LinkedTetraChain is the primary structure for sequence data, DOGMA
also supports general 3D point clouds through Delaunay tetrahedralization.

Given a set of points $P \subset \mathbb{R}^3$, the \emph{Delaunay
tetrahedralization} partitions the convex hull of $P$ into tetrahedra such that
no point lies inside the circumsphere of any tetrahedron.

\begin{definition}[Delaunay Tetrahedralization]
A tetrahedralization of point set $P$ is \emph{Delaunay} if, for every
tetrahedron $\mathcal{T}$, the open circumsphere of $\mathcal{T}$ contains no
point of $P$.
\end{definition}

The Delaunay condition maximizes the minimum solid angle among all
tetrahedralizations, producing well-shaped elements that avoid degeneracies.
This quality guarantee is important for numerical stability in TetraMemory
computations.

Applications relevant to DOGMA:
\begin{itemize}[leftmargin=1.8em]
  \item \textbf{Protein structure:} Given $C_\alpha$ coordinates, Delaunay
        tetrahedralization produces a volumetric mesh capturing local packing
        geometry.
  \item \textbf{Medical imaging:} CT/MRI voxels can be converted to Delaunay
        tetrahedral meshes for volumetric analysis.
  \item \textbf{Molecular dynamics:} Voronoi tessellation (dual of Delaunay)
        defines natural atomic neighborhoods.
\end{itemize}

% ═══════════════════════════════════════════════════════════════════════════════
\section{Multimodal Applications of TDS}
\label{sec:tds-multimodal}
% ═══════════════════════════════════════════════════════════════════════════════

The TDS representation enables DOGMA to handle diverse data modalities through
a unified geometric framework:

\begin{center}
\renewcommand{\arraystretch}{1.15}
\begin{tabularx}{0.95\textwidth}{L{2.5cm} L{3.5cm} Y}
\toprule
\textbf{Modality} & \textbf{TDS Representation} & \textbf{Application} \\
\midrule
DNA/RNA & LinkedTetraChain & Sequence modeling, variant calling \\
3D Vision & Volumetric scene $\to$ tetra mesh & Object recognition, scene understanding \\
Proteins & $C_\alpha$ coordinates $\to$ Delaunay & Structure prediction, drug design \\
Medical & CT/MRI voxels $\to$ tetra mesh & Segmentation, anomaly detection \\
Game engines & Physics bodies $\to$ deformable tetra & Destruction simulation, soft bodies \\
Video & Temporal tetrahedra & Motion tracking, action recognition \\
Audio & Waveform helix encoding & Speaker recognition, music analysis \\
\bottomrule
\end{tabularx}
\end{center}

% ═══════════════════════════════════════════════════════════════════════════════
\section{Worked Examples}
\label{sec:ch10-worked-examples}
% ═══════════════════════════════════════════════════════════════════════════════

\begin{example}[Building a LinkedTetraChain from a DNA Sequence]
\label{ex:build-chain}
Given the sequence $5'$-\texttt{CGTACG}-$3'$ (6 bases), construct the
LinkedTetraChain.

\textbf{Step 1: Assign vertices.}
\begin{center}
\begin{tabular}{ccccccc}
\toprule
Vertex & $v_0$ & $v_1$ & $v_2$ & $v_3$ & $v_4$ & $v_5$ \\
\midrule
Base & C & G & T & A & C & G \\
\bottomrule
\end{tabular}
\end{center}

\textbf{Step 2: Enumerate tetrahedra.}
The chain has $6 - 3 = 3$ tetrahedra:
\begin{align*}
  \mathcal{T}_1 &= (v_0, v_1, v_2, v_3) = (\textsf{C}, \textsf{G}, \textsf{T}, \textsf{A}), \\
  \mathcal{T}_2 &= (v_1, v_2, v_3, v_4) = (\textsf{G}, \textsf{T}, \textsf{A}, \textsf{C}), \\
  \mathcal{T}_3 &= (v_2, v_3, v_4, v_5) = (\textsf{T}, \textsf{A}, \textsf{C}, \textsf{G}).
\end{align*}

\textbf{Step 3: Count combinatorics.}
\begin{itemize}[leftmargin=1.5em]
  \item Vertices: 6
  \item Edges: $3(3) + 3 = 12$
  \item Shared faces: $(v_1,v_2,v_3)$ between $\mathcal{T}_1$ and $\mathcal{T}_2$;
        $(v_2,v_3,v_4)$ between $\mathcal{T}_2$ and $\mathcal{T}_3$.
\end{itemize}

\textbf{Step 4: Compute volume and chirality for $\mathcal{T}_1$.}
Suppose we embed the vertices at:
$\mathbf{v}_0 = (0,0,0)$, $\mathbf{v}_1 = (1,0,0)$,
$\mathbf{v}_2 = (0.5, 0.87, 0)$, $\mathbf{v}_3 = (0.5, 0.29, 0.82)$.

Then:
\[
  V_{\text{signed}} = \frac{1}{6} \det \begin{pmatrix}
    1 & 0.5 & 0.5 \\
    0 & 0.87 & 0.29 \\
    0 & 0 & 0.82
  \end{pmatrix}
  = \frac{1}{6}(1 \cdot 0.87 \cdot 0.82) = \frac{0.7134}{6} \approx 0.119.
\]
Since $V_{\text{signed}} > 0$, the chirality is $\chi = +1$ (right-handed).
\end{example}

\begin{example}[TetraMemory vs.\ Self-Attention Cost]
\label{ex:cost-comparison}
Consider a genomic sequence of length $T = 50{,}000$ with embedding dimension
$d = 512$.

\textbf{Self-attention cost per layer:}
\[
  O(T^2 d) = 50{,}000^2 \times 512 = 1.28 \times 10^{12} \text{ FLOPs}.
\]

\textbf{TetraMemory cost per layer} (with $k = 4$ neighbors):
\[
  O(T \cdot k \cdot d) = 50{,}000 \times 4 \times 512 = 1.024 \times 10^{8} \text{ FLOPs}.
\]

\textbf{Speedup per layer:} $\frac{1.28 \times 10^{12}}{1.024 \times 10^{8}} \approx 12{,}500\times$.

Even accounting for the need for $O(T^{1/3}) \approx 37$ layers for global
propagation in a 3D mesh, the total TetraMemory cost is:
\[
  37 \times 1.024 \times 10^{8} = 3.79 \times 10^{9} \text{ FLOPs},
\]
which is still $\sim$$338\times$ faster than a single self-attention layer.
\end{example}

% ═══════════════════════════════════════════════════════════════════════════════
\section{Exercises}
\label{sec:ch10-exercises}
% ═══════════════════════════════════════════════════════════════════════════════

\begin{enumerate}[label=\textbf{10.\arabic*.}, leftmargin=2.5em]

\item \textbf{[Fundamentals]} Compute the volume and chirality of the tetrahedron
with vertices $(0,0,0)$, $(1,0,0)$, $(0,1,0)$, $(0,0,1)$. Verify that
reflecting one vertex across a coordinate plane reverses the chirality.

\item \textbf{[Construction]} Build the LinkedTetraChain for the DNA sequence
$5'$-\texttt{GATTACA}-$3'$. List all tetrahedra, their base assignments,
and the shared faces between consecutive tetrahedra. How many total edges
does the chain have?

\item \textbf{[Analysis]} For a tetrahedral mesh with $M$ elements, each with at
most 4 face-sharing neighbors, compare the computational cost of one TetraMemory
layer ($O(4Md)$) with one self-attention layer ($O(M^2 d)$). For what value of
$M$ does TetraMemory become 100$\times$ faster?

\item \textbf{[Information Density]} Compute the number of unique pairwise
interactions captured by (a) a triangle mesh with sliding window of width 3
and overlap 2, and (b) a tetrahedral mesh with sliding window of width 4 and
overlap 3, both applied to a sequence of length $T = 100$. Express the ratio.

\item \textbf{[DOF Counting]} A hexahedron (cube) has 8 vertices, 12 edges, and
6 faces. Compute the number of degrees of freedom for a ``HexaData'' structure
analogous to TDS (vertex positions + edge weights + base encoding for 8 bases +
volume + chirality). Is the DOF-per-vertex ratio better or worse than TDS?

\item \textbf{[Coding]} Implement the TetraTokenizer (Equation~\ref{eq:tetra-token})
in Python. Given a list of tetrahedra (each specified by 4 vertex coordinates
and 4 base types), produce the 36-dimensional token representation for each.

\item \textbf{[DoubleHelix]} Given the sense strand $5'$-\texttt{ATGCATGC}-$3'$,
write out the antisense strand. Construct the DoubleHelixMesh and count the
total number of cross-strand hydrogen bonds. How many are A--T bonds (2H) and
how many are G--C bonds (3H)?

\item \textbf{[Theory]} Prove that the graph diameter of a Delaunay
tetrahedralization of $M$ uniformly distributed points in a unit cube is
$O(M^{1/3})$.

\item \textbf{[Design]} Propose a scheme for encoding a protein backbone
(sequence of $C_\alpha$ coordinates) as a TDS mesh. How would you assign base
types $\{\mathsf{A}, \mathsf{C}, \mathsf{G}, \mathsf{T}\}$ to vertices based
on amino acid properties?

\item \textbf{[Research]} Compare TetraMemory with graph neural networks (GNNs).
Both use local message passing. What does the tetrahedral structure provide that
a generic graph does not? Consider: bounded degree, volume, chirality, and
face-sharing.

\end{enumerate}

% ═══════════════════════════════════════════════════════════════════════════════
\section{Key Takeaways}

\keytakeaway{
\begin{itemize}[leftmargin=1.5em]
  \item A tetrahedron is the \emph{minimal solid}: four vertices, six edges, four
        faces, and one volume. It is the simplest 3D object with both volume and
        chirality. No simpler primitive can encode these properties.
  \item The TetraData Structure uses tetrahedra as computational primitives,
        mapping four vertices to DNA base types $\{\mathsf{A}, \mathsf{C},
        \mathsf{G}, \mathsf{T}\}$ and encoding 36 degrees of freedom per element.
  \item The LinkedTetraChain links tetrahedra by shared triangular faces, acting
        as a sliding window of width 4 with overlap 3 over a sequence. Each new
        vertex adds exactly 3 edges and 1 tetrahedron.
  \item DNA sequences map naturally to LinkedTetraChains: each tetrahedron
        represents a window of 4 consecutive bases, and adjacent tetrahedra
        share 3 bases through their common face.
  \item TetraMemory replaces $O(T^2)$ self-attention with $O(T \cdot k)$ local
        tetrahedral interactions ($k \leq 4$) that compose into global computation
        through face propagation. Volume serves as a geometric attention weight;
        edge weights encode local pairwise compatibility.
  \item The DoubleHelixMesh pairs a sense and antisense LinkedTetraChain with
        Watson--Crick cross-links (A--T: 2 hydrogen bonds, G--C: 3 hydrogen
        bonds), providing redundancy, bidirectional reading, and cross-strand
        attention.
  \item Delaunay tetrahedralization extends TDS beyond sequences to general 3D
        point clouds, enabling multimodal applications from protein structure to
        medical imaging.
\end{itemize}
}

\section*{Suggested Reading}
\begin{itemize}[leftmargin=1.5em]
  \item Chapter~\ref{ch:dogma-architecture}, Section on DNAComputeBlock:
        TetraMemory's architectural context.
  \item \citet{jumper2021}: AlphaFold's use of structural representations for
        proteins.
  \item \citet{gu2023mamba}: Alternative approach to sub-quadratic sequence
        processing.
\end{itemize}
